\section{Discussion}\label{sec:discussion}

\paragraph{What the evidence says about learning}
Across the original loop, a standard VND, a ten-operator pool, tuned policies, further policy families (including one that learns the gain per distance evaluation) and the move order of HGS-CVRP (Studies 1--4, 6 and 7), no adaptive policy outperformed uniform random selection (none of \num{u_n} confirmatory contrasts shows a significant improvement, whereas \num{u_worse} show the policy significantly worse), and the policy under study, HG-UCB, was significantly worse than chance in \num{u_hg_worse} of its \num{u_hg_n} contrasts. The equivalence audit rules out gains of adaptation above about \num{eq_best_lo90}~pp, and a static operator mixture tuned offline gains at most \num{s4a_V4_1}~pp over uniform selection at the smaller budget and nothing at the larger one, so there is little headroom for a policy to capture by learning a global preference among these operators. The size of the bound is of the order of the average gain of the adaptive layer in a meta-analysis of adaptive large neighborhood search (an improvement of the objective by 0.14\%, \citealp{turkes2021}, which is, up to the normalization by the reference cost, a difference of 0.14~pp in gap), whereas termination, operator order and the operator pool changed the gap by \num{g_struct_lo} to \num{g_struct_hi}~pp, \num{g_struct_x_lo} to \num{g_struct_x_hi} times more.

\paragraph{Why adaptation has so little to gain here}
Four mechanisms are consistent with our data. (i) \emph{Deterministic operators and local optima.} Each operator applies its best improving move, so once all operators have failed on a solution nothing more can be learned or gained; in the original loop \num{diag_ultima}\% of the applications of a step follow its last improvement and \num{diag_otimo}\% follow the point at which all five operators have failed, which also left a policy only the first few applications of each step to influence. (ii) \emph{Weak and sparse signal.} Most applications fail and the successes are small relative improvements, so the learned policies stay close to uniform (HG-UCB uses each of the five operators \num{s2_hg_use_min}--\num{s2_hg_use_max}\% of the time). (iii) \emph{Cost, not success, separates the operators.} In the ten-operator pool the operators differ by a factor of \num{eff_ratio} in gain per distance evaluation (Table~\ref{tab:eff}), but a distribution tuned on 24 instances to exploit this gains little over uniform selection (Study~4), and a patience stopping rule does not significantly improve on stopping at a local optimum (Study~3). (iv) \emph{Cost of failures.} The cost-aware reward of Studies 3 and 4 credits cheap successes but not the cost of failures; the policy of Study~6 that does charge them is the worst of all (\num{s6d_W5_5}~pp at $E_2$ with $C=1$), because expensive operators are worth more than their efficiency suggests (Study~5).

\paragraph{What does improve the results}
Stopping the inner loop at local optima, randomizing the operator order and enlarging the operator pool changed the gap by more than any policy. Study~5 shows that the pool effect has no single source: the three cheap added operators and the two expensive ones each recover a large part of it, the expensive pair more (\num{s5_rec_P7e_rng}\% against \num{s5_rec_P8c_rng}\%), although they are the least efficient per distance evaluation under uniform selection (Table~\ref{tab:eff}); they are what finds improvements at the local optima of the original operators (Section~\ref{sec:s5}). Efficiency measured in that way is therefore not a safe guide for pruning an operator pool, whereas an ablation of the pool itself is. In applications where the operators are cheap and similar, the practical implication is that the simplest option, a VND with random operator order that stops at local optima, is a strong default, and that effort is better spent on the operator pool and on its cost profile than on the selection layer.

\paragraph{Implications for the design and evaluation of learning-guided heuristics}
Four practices follow. First, \emph{include a matched non-adaptive control}: without it, the policies of Study~2 would appear to improve on the classical fixed order by \num{g_vsfixed_abs}~pp, a gain that belongs to randomization. Second, \emph{check the scale of the reward}, since a reward in cost units can turn a nominally adaptive method into one that is worse than chance. Third, \emph{ablate the search structure and the operator pool}, which moved the gap by more than any policy. Fourth, \emph{report equivalence, not only non-significance}, and add a \emph{tuned static mixture} to bound what adaptation can capture.

\paragraph{Relation to the original study}
We do not claim that the published improvement is wrong: our reconstruction differs from the original implementation in aspects that the pseudocode leaves open (acceptance rule, operator order, cooperative parallel search, compiled code), and our baseline agrees with the published one only on the mean of the three sets. We can say that the effect is not reproduced under a faithful-as-possible reconstruction, that the reward as published is scale-dependent, and that the improvement attributed to learning is smaller than the effects of structure in our experiments.

\subsection{Limitations}
\begin{itemize}
\item \emph{Platform.} Most of our searches are a Python VNS with deterministic best-improvement operators and gaps of \num{vns_gap_min}--\num{vns_gap_max}\%, against \num{hgs_n01}--\num{hgs_matched}\% for a hybrid genetic search. Study~7 tests adaptive ordering of the node-pair moves inside HGS-CVRP, with a binary reward, batched index updates and SWAP* left outside the policy's control; because even randomizing that order changes the gap by at most \num{s7_maxabs}~pp, it is a weak test, and adaptive choice among whole neighborhoods or among the genetic operators of a state-of-the-art solver is not tested. We claim no ranking or competitiveness.
\item \emph{Policies and operators.} Ten adaptive policy families and a static mixture, on two operator pools from one family of neighborhoods. Contextual policies were tested only with the previous operator and its outcome as context; richer solution features, stochastic operators and other reward definitions may behave differently, and the new families of Study~6 and the policies of Study~7 were not tuned.
\item \emph{Tuning.} Policies were tuned on a modest grid with five seeds at one structure and budget, with unequal effort across families, and the tuning noise (\num{s4_tune_noise}~pp) is of the order of the gains of tuning; a larger tuning budget or an automatic configuration procedure \citep{lopezibanez2016} might yield a better adaptive policy.
\item \emph{Instances.} 74 CVRPLIB instances with up to 194 customers and 54 synthetic ones with up to 300; synthetic reference costs are the best of two HGS runs, not proven optima. Only the CVRP is studied.
\item \emph{Inference.} Tests use instance means, so seed variance is summarized, not modeled; two-level intervals address this only partially. The equivalence margin (except in Study~7) and the robustness analyses were chosen after the confirmatory tests.
\item \emph{Pre-specification.} Hypotheses were written in a dated research log before each study ran, but without an external registry; version-control timestamps exist for Study~4 (commit \texttt{ef82cdc}), Study~6 (\texttt{02e58d3}, extended by \texttt{b3d4add}) and Study~7 (\texttt{8a06b17}), each before its data; the entry of Study~5 was committed after its runs had started, and Studies 1--3 rest on the log alone.
\item \emph{Reproduction.} The published implementation could not be obtained; the reproduction rests on the published description and four documented interpretation choices.
\end{itemize}

\section{Conclusions}\label{sec:conclusions}
Auditing a published bandit-guided VNS for the CVRP, we found that its reported improvement is not reproduced, that the reward as published can reduce it to worse-than-random behavior, and that, once the reward is normalized, none of the adaptive policies tested beats uniform random selection: not in the original loop, inside a standard VND, with a heterogeneous operator pool, after tuning, nor with further bandit families and a policy that learns the gain per distance evaluation, nor for the move order inside the local search of a state-of-the-art genetic algorithm, a weak test because that order barely matters. The gains came from elsewhere: stopping the inner loop at local optima, randomizing operator order and enlarging the operator pool. Both cheap and expensive added operators contribute to the pool effect, the expensive ones more than their efficiency per evaluation suggests. For operations research the lesson is methodological as much as algorithmic: the four practices listed above should precede any claim about the contribution of a learned component. Natural extensions are testing adaptation at other levels of state-of-the-art solvers, richer contextual policies, and other routing problems; our protocol, data and code are provided to support them.
